% Ch5 WS3 -- KMT, real gases, plus an FRQ. Block 5.
\documentclass{apchem-worksheet}

\apcunitword{Chapter}
\apcunit{5}
\apcunittitle{Gases}
\apcdoctitle{Worksheet 3 \textbullet\ KMT, Real Gases \& FRQ}
\apcsubtitle{Zumdahl \S5.6--5.9 \textbullet\ name the failed postulate, not just ``it's real''}

\begin{document}
\wsheader[Deviation answers must name \emph{which} postulate fails and
\emph{which direction} $PV/nRT$ moves. ``Because it is a real gas'' earns
nothing.]

\problem[]{State which KMT postulate each observation relies on:}
\begin{parts}
  \item A gas fills its container uniformly.
        \answerblank[5.5cm]{constant random motion}
  \item Gas pressure in a sealed flask does not decay over time.
        \answerblank[5.5cm]{elastic collisions}
  \item \ce{He} and \ce{Xe} at the same $T$ have equal average kinetic
        energies. \answerblank[5.5cm]{$KE \propto T$}
\end{parts}

\problem[]{At \SI{300}{\kelvin}, rank \ce{H2}, \ce{O2}, and \ce{SF6}
($M = 146.06$) by:}
\begin{parts}
  \item Average kinetic energy:
        \answerblank[6.3cm]{all equal --- depends only on $T$}
  \item Average molecular speed:
        \answerblank[5cm]{\ce{H2} $>$ \ce{O2} $>$ \ce{SF6}}
  \item Explain how both answers can be true at once.
        \answerlines[2]{$KE = \tfrac12 mu^2$: heavier molecules compensate
        for greater mass by moving more slowly, so all reach the same
        average energy at a given temperature.}
\end{parts}

\problem[]{Sketch, on one set of axes, Maxwell--Boltzmann distributions for
one gas at \SI{200}{\kelvin} and at \SI{600}{\kelvin}. Label both curves and
state what stays the same. \workspace[4.2cm]
\keyonly{\begin{apcanswer}600 K curve: broader, flatter, peak shifted right.
200 K: narrower, taller, peak left. The \emph{area} under both is the same
--- it represents the total number of molecules.\end{apcanswer}}}

\problem[]{Real gas behaviour:}
\begin{parts}
  \item Under what two conditions does a real gas behave most ideally?
        \answerblank[5.5cm]{high $T$, low $P$}
  \item At high pressure, $PV/nRT$ rises above 1. Name the failed postulate
        and explain. \answerlines[2]{``Particles have negligible volume''
        fails --- the molecules' own volume is a significant fraction of the
        container, so the gas cannot be compressed as far as predicted.}
  \item At low temperature, $PV/nRT$ dips below 1. Name the failed postulate
        and explain. \answerlines[2]{``No attractive forces'' fails ---
        molecules pull on one another, softening wall collisions and
        lowering the measured pressure below the ideal value.}
  \item Which deviates more, \ce{Ne} or \ce{NH3}? Why?
        \answerblank[6cm]{\ce{NH3} --- hydrogen bonding}
\end{parts}

\problem[]{\textbf{FRQ (10 points).} A student generates oxygen by
decomposing potassium chlorate and collects it over water:
\[ \ce{2 KClO3(s) -> 2 KCl(s) + 3 O2(g)} \]
A sample of \SI{2.45}{\gram} \ce{KClO3} ($M = \SI{122.55}{\gram\per\mole}$)
is fully decomposed. The gas is collected at \SI{25}{\celsius} with a total
pressure of \SI{754}{\torr}. The vapour pressure of water at
\SI{25}{\celsius} is \SI{23.8}{\torr}.}
\begin{parts}
  \item Calculate the moles of \ce{O2} that should be produced.
        \workspace[2.2cm]
        \keyonly{\begin{apcanswer}$n(\ce{KClO3}) = 2.45/122.55 =
        \SI{0.0200}{\mole}$; $\times\tfrac32 =
        \SI{0.0300}{\mole}$ \ce{O2}\end{apcanswer}}
  \item Determine the partial pressure of \ce{O2}, in atm.
        \answerlines[2]{$754 - 23.8 = \SI{730.2}{\torr}$;
        $730.2/760 = \SI{0.9608}{\atm}$}
  \item Calculate the volume of gas collected. \workspace[2.2cm]
        \keyonly{\begin{apcanswer}$V = nRT/P = (0.0300)(0.08206)(298)/0.9608
        = \SI{0.764}{\liter}$\end{apcanswer}}
  \item Explain why the water-vapour correction is a subtraction, and state
        the effect of omitting it. \answerlines[3]{The collected gas is a
        mixture, so the measured total includes water vapour's contribution;
        the oxygen alone accounts for less. Omitting the correction
        overstates $P_{\ce{O2}}$ and therefore understates the calculated
        volume for a known number of moles.}
  \item The student repeats the collection at \SI{35}{\celsius}
        ($P_{\ce{H2O}} = \SI{42.2}{\torr}$), same total pressure. Predict
        whether the collected volume is larger or smaller, and give both
        reasons. \answerlines[3]{Larger. Higher temperature increases volume
        directly ($V \propto T$), and the larger water-vapour correction
        lowers $P_{\ce{O2}}$, which also increases the volume needed to hold
        the same moles.}
\end{parts}

\begin{apcnote}[Rubric --- score yourself, 10 points]
\textbf{(a) 2 pts} 1 moles of \ce{KClO3}, 1 the 3:2 ratio applied
\textbullet\ \textbf{(b) 2 pts} 1 subtraction performed, 1 conversion to atm
\textbullet\ \textbf{(c) 2 pts} 1 correct rearrangement with kelvin, 1 value
with units \textbullet\ \textbf{(d) 2 pts} 1 explains the mixture/Dalton
reasoning, 1 states the direction of the error \textbullet\
\textbf{(e) 2 pts} 1 predicts larger, 1 gives BOTH reasons --- one reason
only earns 1.
\end{apcnote}

\begin{spiralstrip}{Chapter 5 \textbullet\ blocks 1--4}
  \item $P_{\ce{H2O}}$ at \SI{30}{\celsius}: \answerblank[2.6cm]{\SI{31.8}{\torr}}
  \item Volume of \SI{3.00}{\mole} at STP: \answerblank[2.8cm]{\SI{67.3}{\liter}}
  \item $\chi_A$ if $P_A = \SI{1.5}{\atm}$, $P_{\text{tot}} =
        \SI{6.0}{\atm}$: \answerblank[2cm]{0.25}
  \item Faster at equal $T$: \ce{CH4} or \ce{Ar}?
        \answerblank[2.4cm]{\ce{CH4}}
  \item $d$ of a gas with $M = \SI{28.0}{\gram\per\mole}$ at STP:
        \answerblank[3.2cm]{\SI{1.25}{\gram\per\liter}}
\end{spiralstrip}

\end{document}
